% Related Work — Article 1
% "Attestation-Aware Scheduling for Verifiable AI Placement on SEV-SNP
%  Confidential Kubernetes Nodes"
%
% Every \cite key below resolves in references_verified.bib (all real, verified).
% This section is written to pre-empt the Q1 reviewer objections tracked in
% article1/status/reviewer_attack_matrix.md.

\section{Related Work}
\label{sec:related}

We position our contribution against six strands of prior work: (i) TEE hardware
and confidential VMs, (ii) remote attestation architecture and protocols,
(iii) TEE systems and secure containers, (iv) confidential containers and
confidential Kubernetes, (v) confidential ML/AI inference, and (vi) GPU trusted
execution. Throughout, we make explicit what we do \emph{not} claim: our
evaluation targets AMD SEV-SNP \emph{node-level} attestation on Confidential VM
node pools, not per-pod TEE attestation and not confidential GPUs.

\subsection{TEE hardware and confidential VMs}
Confidential Virtual Machines shift the trust boundary from a process-scoped
enclave to a whole encrypted VM. AMD SEV-SNP and Intel TDX are the dominant CVM
substrates; their architectures, interfaces and security properties differ in
subtle ways that complicate reasoning about evidence semantics
\cite{misono2024cvms,cheng2023tdx}. Formal analysis of the SEV-SNP software
interface has proved most secrecy, authentication, attestation and freshness
properties in a symbolic model while flagging residual weaknesses
\cite{paradzik2024sevsnp}, and micro-architectural side channels against SEV
(e.g. the ciphertext side channel \cite{li2021cipherleaks}) motivate treating a
confidential node's evidence as necessary but bounded --- reinforcing our
freshness and revocation checks rather than assuming a node is trusted forever. Foundational enclave designs (Intel SGX
\cite{costan2016sgx}, the open RISC-V Keystone framework \cite{lee2020keystone})
and Arm's confidential-compute direction \cite{xu2023virtcca,cerdeira2022rezone}
establish the primitives our evidence is ultimately rooted in. Recent
systematizations argue that a missing \emph{abstraction layer} hinders broad CVM
adoption \cite{michaud2025teesok}; our work contributes one such layer at the
orchestration tier: verified attestation evidence as a first-class scheduling
signal. Crucially, none of these works addresses \emph{where} in a cluster a
sensitive workload is placed, nor binds a placement decision to the evidence.

\subsection{Remote attestation architecture and protocols}
The IETF RATS architecture \cite{birkholz2023rats} standardises the
Attester/Verifier/Relying-Party roles and the evidence$\rightarrow$attestation-result
flow. A recurring, security-critical requirement is the strict separation between
the party that \emph{produces} evidence and the party that \emph{appraises} it;
open attestation systems such as OPERA \cite{chen2019opera} and RA-TLS-style
trusted channels \cite{knauth2018ratls,akama2024raweb} operationalise this
separation for enclaves, and in-enclave verification of policy compliance has
been explored \cite{liu2020confidattest}. Bridging attestation technologies
(e.g., attesting pre-SNP SEV VMs via SGX enclaves) further underlines that
appraisal must be an independent, trusted step \cite{antonino2023flexattest}.
Our design embodies the RATS separation \emph{inside Kubernetes}: a
node-resident agent may only emit an untrusted \texttt{RawAttestationReport},
while a single central verifier is the sole writer of the appraised
\texttt{AttestationEvidence}. This directly answers the reviewer questions
``who verifies attestation?'' and ``can a node self-attest?'' --- it cannot.

\subsection{TEE systems and secure containers}
Library-OS and secure-container systems (SCONE \cite{arnautov2016scone},
Graphene/Gramine \cite{tsai2017graphene}, Occlum \cite{shen2020occlum}) run
unmodified or lightly modified workloads inside enclaves, and language- and
storage-level confidential runtimes extend the model
\cite{sarkar2024hastee,hartono2024crisp}. These systems harden the \emph{inside}
of a confidential workload. They are orthogonal and complementary to our work,
which does not modify the runtime: we govern \emph{admission and placement} so
that a sensitive pod only ever binds to a node whose attestation evidence has
been independently verified, fresh, and non-revoked.

\subsection{Confidential containers and confidential Kubernetes}
The closest ecosystem is Confidential Containers (CoCo) with the Trustee
attestation/secret-release service \cite{coco2024}, and the broader
``confidential Kubernetes'' pattern of running nodes or pods inside CVMs
\cite{k8s2023confidential}. Pod-level remote attestation for Kubernetes
workloads has recently been demonstrated on dstack \cite{yang2026dstack}, and
empirical work has characterised trust-boundary vulnerabilities in TEE
containers \cite{liu2025teecontainers}. We deliberately do \emph{not} claim
pod-level attestation: on AKS SEV-SNP Confidential VM node pools the evidence
refers to the confidential worker node/VM, so we report \emph{attested-node
placement}. Our contribution is orthogonal to CoCo/Trustee: rather than
gating secret \emph{release} at runtime, we make verified evidence a
\emph{scheduling-time} signal and bind the final \texttt{AIPlacementDecision} to
independently verifiable cryptographic evidence via a signed placement token.
Compared with industrial confidential-Kubernetes distributions (e.g.
Constellation-style whole-cluster CVMs, and C8s-style confidential control
planes), which secure the cluster substrate, we target scheduling-time
enforcement and a verifiable placement artifact; the two are composable.

\paragraph{Why not simpler mechanisms?}
A Q1 reviewer will ask why node labels, \texttt{RuntimeClass}, or
\texttt{schedulingGates} plus an external controller do not already suffice.
Labels and \texttt{nodeSelector} are attacker-forgeable cluster metadata with no
cryptographic backing; \texttt{RuntimeClass} selects a runtime handler but
carries no freshness, revocation, or node-identity semantics; and a
\texttt{schedulingGates}+controller design (our strong baseline B4) can check
evidence \emph{before} ungating but leaves a Time-Of-Check-to-Time-Of-Use
(TOCTOU) window between the check and the bind, and produces no offline-verifiable
placement artifact. We quantify this TOCTOU window against our fail-closed
\texttt{PreBind} re-check and show that only our design emits a signed token that
an independent \texttt{verify-placement} tool validates without cluster access.

\subsection{Confidential ML/AI inference}
A large body of work protects models and inference inside TEEs, via DNN
partitioning between trusted and untrusted compute (Slalom \cite{tramer2019slalom},
DarkneTZ \cite{mo2020darknetz}, ShadowNet \cite{sun2020shadownet},
TEESlice \cite{li2024teeslice}), secure accelerators
\cite{hua2020guardnn,hashemi2022darknight}, and secure serving systems
\cite{ma2020s3ml,lee2020ppml,hu2024sesemi,schambach2026enclavex,yu2025dualprivacy,abdollahi2025ondevice,abdollahi2026agentee}.
Systematizations map the design space and its pitfalls
\cite{mo2022mlccsok,duy2021cmlsok}, and TEE-shielded partitioning has itself been
shown insecure under strong adversaries \cite{zhang2023noprivacy}. These works
answer \emph{how} to run a confidential model; they assume the workload already
runs on a trustworthy node. Our work answers the prior, orchestration-level
question --- \emph{how to guarantee, verifiably, that it was placed on an
attested node} --- and is therefore complementary to all of them.

\subsection{GPU trusted execution (future work)}
GPU TEEs (Graviton \cite{volos2018graviton}, Telekine \cite{hunt2020telekine})
and production confidential GPUs (NVIDIA Hopper H100, analysed for performance
\cite{zhu2024h100bench} and security \cite{gu2025gpucc}) extend confidential
execution to accelerators. We do not evaluate confidential GPUs: the DCasv6
quota used here is CPU-only SEV-SNP. GPU-aware placement is expressed in our
policy schema and left as validated future work with NCC H100 hardware.

\subsection{Orchestration integrity and verifiable provenance}
Finally, supply-chain and provenance work --- hardware-attested ML property
cards \cite{duddu2024laminator}, and hardened Kubernetes deployments
\cite{ali2024cmsweb,enyedi2026modrepos} --- shares our goal of making
security claims \emph{verifiable} rather than asserted. Our signed
\texttt{AIPlacementDecision} extends this ethos to the placement decision itself:
the ``where did my confidential workload run?'' question becomes answerable
offline by a third party.

% IP REVIEW REQUIRED BEFORE SUBMISSION
% (This section references the placement token / signed AIPlacementDecision.)
